\documentclass[11pt]{article}

% Self-contained: compiles on Overleaf or any standard TeX install with no extra class files.
\usepackage[utf8]{inputenc}
\usepackage[T1]{fontenc}
\usepackage{amsmath,amssymb}
\usepackage{microtype}
\usepackage[margin=1in]{geometry}
\usepackage{hyperref}
\hypersetup{colorlinks=true,linkcolor=blue,citecolor=blue,urlcolor=blue}
\usepackage{xcolor}
\usepackage{listings}
\lstset{basicstyle=\ttfamily\small,breaklines=true,columns=fullflexible,frame=single,
        framesep=4pt,xleftmargin=4pt}

\newcommand{\Q}{Q}
\newcommand{\code}[1]{\texttt{#1}}

\title{Where SMT-Based Verification of Post-Quantum Modular Arithmetic Cliffs:\\
       A Cross-Toolchain Tractability Study of ML-DSA}

% TODO(author): confirm name / affiliation / contact before submission.
\author{Amar Akshat\\\texttt{amar.akshat@gmail.com}}
\date{\today}

\begin{document}
\maketitle

\begin{abstract}
Machine-checked equivalence between a cryptographic implementation and its specification is only
useful if the proof actually discharges. For the modular arithmetic at the core of ML-DSA
(FIPS~204)~\cite{fips204} we find that whether it discharges depends sharply, and predictably, on the
input domain --- and we map where the boundary lies, why it is there, and how to cross it. We verify
reduction primitives and the forward number-theoretic transform (NTT) of two independent, deployed
ML-DSA codebases: the PQClean / pq-crystals C reference, via SAW, Cryptol, and Isabelle/HOL composed
through the \code{cryptol-to-isabelle} translator into a single chain from C to a FIPS-derived
specification; and the RustCrypto implementation, via SAW over MIR. On the C side we obtain end-to-end
correctness for the \code{reduce.c} layer (\code{montgomery\_reduce}, \code{reduce32}, \code{caddq},
\code{freeze}) and a machine-checked overflow-freedom bound for the forward-NTT model. On the Rust side
we obtain SAW-verified equivalence for Barrett reduction, constant-time division, the twiddle-factor
table, and the hint layer (Algorithms 36, 37, 39, 40). Against these positive results we report the
boundary: narrow-domain reductions discharge in seconds, while wide-domain Barrett reduction over the
full product domain $x < 2^{46}$, and a single NTT layer, fail to complete across six solver and tactic
configurations. We give the structural cause --- bit-blasting $(x \cdot M) \gg k$ over a wide symbolic
domain --- and three constructive escapes that do discharge: a total-coefficient-invariant induction
that replaces a $\sim$3000-obligation SAW unrolling with an eight-step proof running in seconds, an
unbounded-integer reformulation with a bitvector bridge, and compositional overrides. We also report
two off-by-one defects in the reference's documented contracts, disclosed upstream. The development uses
no vendor verification artifacts and is reproducible with a single command.
\end{abstract}

\medskip
\noindent\textbf{Keywords:} ML-DSA $\cdot$ FIPS 204 $\cdot$ formal verification $\cdot$ SMT
$\cdot$ SAW $\cdot$ Cryptol $\cdot$ Isabelle/HOL $\cdot$ MIR $\cdot$ Barrett reduction $\cdot$
number-theoretic transform $\cdot$ tractability.

\section{Introduction}\label{sec:intro}

ML-DSA (FIPS~204)~\cite{fips204} is the NIST-standardized lattice signature scheme derived from
CRYSTALS-Dilithium~\cite{dilithium}. Its security analysis assumes that implementations compute the
specified function. Test vectors check that assumption on a finite set of inputs; formal verification
checks it for all inputs, which matters for arithmetic kernels whose bugs are input-dependent and can
pass test vectors~\cite{eprint1032}.

A practitioner who sets out to machine-check these kernels against the specification quickly meets a
less-discussed obstacle: many of the proofs do not discharge. The same SAW/SMT formulation that proves
one reduction in seconds runs for hours, or never finishes, on another that differs only in its input
range. This is folklore among verification engineers, but it is rarely characterized: \emph{where} the
boundary sits, \emph{why}, and \emph{what to do instead} are usually left implicit.

This paper maps that boundary for ML-DSA's modular arithmetic, concretely and reproducibly, across two
deployed implementations and two verification toolchains. We treat the cases that do not discharge not
as failures to hide but as the primary result: a characterized tractability boundary, its structural
cause, and the proof restructurings that cross it. The positive verification results --- a complete
reduction layer and an overflow-freedom bound on the C side, and a verified field, division,
twiddle-table, and hint layer on the Rust side --- establish what sits on the tractable side of the
boundary and anchor the study.

\paragraph{Contributions.}
\begin{enumerate}
  \item \textbf{A reproducible tractability map} for SMT-based equivalence verification of PQC modular
  arithmetic (Section~\ref{sec:boundary}). We characterize, with timings, which arithmetic-kernel
  equivalences discharge automatically and which do not: narrow-domain Barrett reduction ($x < 2^{28}$)
  proves in about four seconds, while the same goal over the cryptographically relevant wide domain
  ($x < 2^{46}$) does not complete under \code{z3}, \code{abc}, \code{cvc5}, \code{yices}, a
  \code{goal\_eval} rewrite pass, or a carry-fold reformulation.
  \item \textbf{The structural cause.} The wall is bit-blasting of $(x \cdot M) \gg k$ --- a wide
  product divided by a power of two --- over a wide symbolic input domain, the shape that recurs in
  Barrett reduction and, 256-fold, in one NTT layer (Section~\ref{sec:boundary}). This predicts the
  boundary rather than merely observing it.
  \item \textbf{Three constructive escapes that discharge.} A total-coefficient invariant turns an
  intractable $\sim$3000-obligation SAW unrolling of NTT overflow-freedom into an eight-step Isabelle
  induction running in seconds (Section~\ref{sec:ntt}); an unbounded-integer reformulation with a
  bitvector bridge handles the reductions; compositional overrides handle the layer proofs
  (Sections~\ref{sec:reduce} and~\ref{sec:boundary}).
  \item \textbf{A cross-toolchain, cross-language study.} The same primitives are verified in two
  pipelines --- SAW+Cryptol+Isabelle on the C reference (Section~\ref{sec:reduce}) and SAW over MIR on
  the RustCrypto implementation (Section~\ref{sec:rust}). We are not aware of a prior side-by-side
  machine verification of ML-DSA reduction code across two toolchains and two languages; it lets us show
  the boundary is a property of the goal shape, not of one toolchain.
  \item \textbf{Two off-by-one defects} in the reference's documentation contracts, reproduced against
  the shipping code and disclosed upstream to pq-crystals/dilithium (Section~\ref{sec:findings}).
\end{enumerate}

\paragraph{Non-claims.} We verify reference and library implementations, not the optimized
AVX2/aarch64 paths or the \code{mldsa-native} successor, where the most error-prone arithmetic lives. We
prove C/Rust $\equiv$ model and overflow-freedom; we do not prove NTT \emph{transform} correctness, that
is, equivalence of the model to the negacyclic transform of FIPS~204 Algorithms 41 and 42. The
wide-domain reductions and the NTT-layer transform proof are reported as a characterized boundary with
escapes identified, \emph{not} as completed proofs. We say nothing about constant-time behavior.

\section{Background}\label{sec:background}

\paragraph{Targets.} On the C side we verify PQClean~\cite{pqclean} \code{ml-dsa-44/clean}, files
\code{reduce.c} and \code{ntt.c}, pinned by commit hash; this code is byte-identical to
pq-crystals/dilithium~\cite{dilithium} up to a namespace prefix. On the Rust side we verify
RustCrypto's \code{ml-dsa} crate~\cite{rustcrypto} and its shared \code{module-lattice} field core. ML-DSA
fixes the prime modulus $Q = 8380417$ and ring dimension $N = 256$. \code{montgomery\_reduce} maps a
64-bit input to a 32-bit value congruent to $a \cdot 2^{-32} \pmod Q$. \code{reduce32} is a
Barrett-style reduction to a centered window; RustCrypto instead uses an explicit Barrett reduction
$x \mapsto x - \lfloor (x \cdot M) / 2^{k} \rfloor\, Q$ with $M = \lfloor 2^{k}/Q \rfloor$. The forward
NTT applies eight levels of Cooley-Tukey butterflies, each multiplying a coefficient by a precomputed
twiddle factor and reducing.

\paragraph{Tools.} SAW (the Software Analysis Workbench)~\cite{saw} proves a compiled function
equivalent to a Cryptol specification by symbolic execution and SMT; it targets LLVM bitcode for C and,
via \code{mir-json}, Rust MIR. Cryptol is a word-level specification language. Isabelle/HOL is an
interactive theorem prover. The \code{cryptol-to-isabelle} translator, first shipped with saw-script
1.5.1, lifts a Cryptol module into an Isabelle theory; we use it to compose the two C-side proof legs.
All tool versions are pinned (Section~\ref{sec:repro}).

\section{The two pipelines}\label{sec:pipeline}

\begin{figure}[h]
\centering
\begin{lstlisting}[frame=none]
  C reference                         RustCrypto (Rust)
      |  hand-translate                    |  mir-json
      v                                     v
  Cryptol model --SAW: C == Cryptol-->   MIR  --SAW: Rust == Cryptol-->  proven
      |  cryptol-to-isabelle
      v
  Isabelle model --model == FIPS spec--> proven (reduce.c layer)
\end{lstlisting}
\caption{Two pipelines over the same primitives. The C pipeline composes a SAW leg and an Isabelle leg
through \code{cryptol-to-isabelle}, giving a chain from C to a FIPS-derived specification. The Rust
pipeline verifies MIR directly against Cryptol with SAW. Both meet the same tractability boundary
(Section~\ref{sec:boundary}).}
\label{fig:pipeline}
\end{figure}

The C model is hand-translated at the exact bit widths, because the algorithms depend on two's-complement
truncation and signed shifts and SAW checks bit-for-bit equivalence. The lift step is mechanical: a
continuous-integration check (\code{make lift-check}) regenerates the Isabelle theory from the
\code{.cry} source and diffs it against the committed copy, so no human-maintained link sits between the
two legs. The Rust pipeline builds MIR with \code{mir-json}~\cite{mirjson} (pinned to the schema version
SAW~1.5.1 bundles) from a small harness crate that exercises keygen, sign, and verify so the generic
field and transform code monomorphizes; SAW then verifies the monomorphized functions against Cryptol.

\section{The \code{reduce.c} layer (C)}\label{sec:reduce}

\paragraph{C $\equiv$ Cryptol.} SAW discharges bit-exact equivalence for all four functions on the
default (overflow-checking) LLVM bitcode, so the proofs also assert absence of signed-overflow undefined
behavior on the relevant ranges. \code{montgomery\_reduce} is verified under its documented input range.
\code{reduce32} is verified under $a \le 2^{31} - 2^{22} - 1$, the bound that keeps the internal
\code{a + (1{<}{<}22)} addition from overflowing. \code{caddq} is verified unconditionally. \code{freeze}
is verified compositionally using the \code{reduce32} and \code{caddq} results as overrides. A mutation
test confirms the \code{montgomery\_reduce} proof is non-vacuous: a deliberately wrong model is rejected
with a counterexample.

\paragraph{Model $\equiv$ specification.} In Isabelle we state each function's contract and prove the
lifted model meets it, with no \code{sorry} or \code{oops}.
\begin{itemize}
  \item \code{montgomery\_reduce}: $2^{32} r \equiv a \pmod Q$ and $-Q < r < Q$ on the half-open domain
  $-2^{31} Q \le a < 2^{31} Q$.
  \item \code{caddq}: residue-preserving, and maps $[-Q, Q)$ into $[0, Q)$.
  \item \code{reduce32}: residue-preserving, with output in the window $[-6283009, 6283008]$ on
  $a \le 2^{31} - 2^{22} - 1$.
  \item \code{freeze}: residue-preserving, with output in $[0, Q)$, proven from the previous two.
\end{itemize}
Composing with the SAW leg gives a statement from C to specification: on the documented ranges, each
deployed reference function computes a correct residue modulo $Q$ within its proven output window. The
windows above are the \emph{true reachable} windows, which differ from the documented ones for two
functions (Section~\ref{sec:findings}). This layer sits entirely on the tractable side of the boundary:
the inputs are single machine words and the goals discharge in seconds.

\section{The RustCrypto layer (Rust)}\label{sec:rust}

The RustCrypto pipeline reaches the same primitives in a different language and toolchain, and the same
narrow-domain goals discharge. All results below are SAW-verified against a Cryptol specification and
checked for non-vacuity by a mutation test (a deliberately altered modulus or constant is rejected).
\begin{itemize}
  \item \textbf{Barrett reduction} $== x \bmod M$ for all \code{u32} inputs, for each of the three moduli
  that occur ($Q = 8380417$, $2\gamma_2 = 190464$, and $2^d = 8192$), with the conditional final
  subtraction live; no precondition is needed because each modulus exceeds $2^{16}$.
  \item \textbf{Constant-time division} $\code{ct\_div}(x) == \lfloor x / 190464 \rfloor$ on its
  documented domain $x < Q$.
  \item \textbf{Twiddle-factor table}: every constant-folded copy of the 256-entry table equals
  $\zeta^{\mathrm{BitRev}_8(i)} \bmod Q$.
  \item \textbf{Hint layer}: \code{decompose}, \code{high\_bits}, \code{make\_hint}, and
  \code{use\_hint} equal FIPS~204 Algorithms 36, 37, 39, and 40 respectively, over all field elements.
\end{itemize}
Three leaf operations are constant-time intrinsics or inline assembly (a \code{black\_box} identity and
two conditional-move primitives); these are supplied as assumed specifications, each justified against
the target assembly, and recorded as the verification's trust base alongside \code{mir-json} soundness.

\section{Forward-NTT overflow-freedom, and the first escape}\label{sec:ntt}

The forward NTT performs unreduced \code{int32} additions and subtractions, $a[j] \pm t$, which overflow
for unbounded inputs. On a \code{-fwrapv} build, where signed overflow is two's-complement wrapping, SAW
proves the C NTT functionally equivalent to the Cryptol model for all inputs, with \code{montgomery\_reduce}
supplied as an uninterpreted override so the 1024 butterfly calls collapse to a structural array
equality. This proof says nothing about overflow.

Overflow-freedom is the property of interest, since it is the defect class behind known ML-DSA
implementation bugs that passed test vectors~\cite{eprint1032}. The statement \code{ntt\_overflow\_free}
is: if every input coefficient lies in $\pm(2^{31} - 2^{27})$, then every coefficient of the transform
lies in $\pm 2080309256 < 2^{31} - 1$. We prove it by induction over the eight levels. A single per-level
lemma states that one level grows the maximum coefficient magnitude by at most $Q$; the growth bound
comes from the \code{montgomery\_reduce} output bound $|t| < Q$. Iterated eight times from an input bound
with headroom $2^{27} > 8Q$, every intermediate value stays below $2^{31}$, so each \code{int32}
addition and subtraction has its true result in range and does not overflow.

\paragraph{The total invariant (escape one).} The enabling device is the form of the invariant. The
Cryptol index operation is strict, but an out-of-bounds access returns the last element of the sequence.
We therefore state the coefficient bound over \emph{all} indices, not only the valid ones. Every indexed
access is then a bounded list member, in-bounds or not, so the proof never reasons about the modular
index arithmetic of the butterfly network --- shifts and divisions parameterized by the level, the exact
arithmetic that made the direct SMT approach intractable (Section~\ref{sec:boundary}). This is the first
of the three escapes: a change of proof structure, not of solver, turns an intractable goal tractable.

\paragraph{From the model to the C.} The Isabelle result is about the model. SAW proves the C equal to
the model for all inputs on the \code{-fwrapv} build; the Isabelle result shows the model does not wrap
under the input bound; the two builds differ only in whether signed overflow is poison. Hence the
reference C NTT has no signed-overflow undefined behavior under the bound. We state this composition
explicitly rather than mechanize it as one SAW theorem, because the direct attempt did not complete
(Section~\ref{sec:boundary}).

\section{The tractability boundary}\label{sec:boundary}

The reduction layer and the field, division, and hint layers above all discharge in seconds. The
transform layer does not, and the failure is sharp and reproducible.

\paragraph{Wide-domain Barrett.} RustCrypto's Barrett reduction computes, in \code{u128},
$\code{product} = x \cdot M$ with $M = \lfloor 2^{46}/Q \rfloor = 8396807$, $\code{quotient} =
\code{product} \gg 46$, $\code{remainder} = x - \code{quotient} \cdot Q$, then a conditional subtraction.
The field multiplication that feeds it produces $x = a \cdot b$ with $a, b < Q < 2^{23}$, so the relevant
input domain is $x < 2^{46}$. Verifying $\code{barrett\_reduce}(x) == x \bmod Q$ against this domain
scales as follows:
\begin{center}
\begin{tabular}{ll}
$x < 2^{20}$ & proves, $\sim$4\,s (z3) \\
$x < 2^{28}$ & proves, $\sim$4\,s (z3) \\
$x < 2^{36}$ & no completion \\
$x < 2^{44}$ & no completion \\
$x < 2^{46}$ & no completion \\
\end{tabular}
\end{center}
At $x < 2^{46}$ the goal did not complete under any of six configurations: \code{z3}, \code{abc},
\code{cvc5}, \code{yices}, a \code{goal\_eval} rewrite pass before \code{z3}, and a hand-written
carry-fold reformulation of the specification (each run to a multi-minute timeout). The corresponding
narrow-domain Barrett reduction over \code{u32} inputs, used elsewhere in ML-DSA, proves without trouble
(Section~\ref{sec:rust}); the difference is the domain width alone.

\paragraph{The structural cause.} The decisive subterm is $(x \cdot M) \gg 46$: a wide product, here up
to roughly $2^{69}$, divided by a power of two, over a symbolic input ranging across $2^{46}$ values.
Bit-blasting this division-of-a-wide-product to a SAT instance is exponential in the operative width, and
no choice of back-end solver changes the shape of the instance. The boundary is therefore a property of
the goal, and it is predictable: it appears wherever a wide-domain $(x \cdot M) \gg k$ must be related to
$x \bmod Q$ at bit-vector width.

\paragraph{The NTT layer, same cause $\times 256$.} The same shape recurs in the transform. We attempted
one forward NTT layer (\code{ntt\_layer} with $\mathrm{LEN}=128$) against a FIPS~204 Algorithm~41 layer
specification, with the four field operations supplied as assumed specifications via
\code{mir\_unsafe\_assume\_spec} so the butterfly reduces to specification-form arithmetic. The overrides
apply and symbolic simulation completes, but the final proof obligation does not discharge: the
specification computes $(\cdot) \bmod Q$ independently for each of 256 output coefficients ---
256 bit-vector remainder-by-constant operations --- and these do not term-match the implementation's
reduce-after-each-operation structure, so the solver bit-blasts all of them. This is the wide-modular
reduction wall again, multiplied across the array.

\paragraph{Escapes two and three.} The total-invariant induction of Section~\ref{sec:ntt} is the first
escape and the one we carried to completion. Two others are available and are the route past the
remaining walls. (2) An \emph{unbounded-integer reformulation}: prove the Barrett identity
$x - ((x \cdot M)\gg k)\,Q \equiv x \pmod Q$ and $\in [0, 2Q)$ in unbounded-integer arithmetic, which a
solver discharges by nonlinear reasoning rather than bit-blasting, then bridge bit-vector to integer ---
the structure by which we proved \code{montgomery\_reduce} on the C side. (3) A \emph{compositional
specification}: build the layer specification from the same override terms the implementation uses, so
the two sides are syntactically the same term and the obligation discharges by rewriting rather than
bit-blasting 256 independent remainders. We report these as the identified escapes; mechanizing the
transform layer through them is the immediate next step.

\section{Findings}\label{sec:findings}

Establishing the true output windows in Section~\ref{sec:reduce} surfaced two off-by-one errors in the
reference's documentation comments. Both are documentation-contract defects, not miscomputations: the
functions return correct residues, but the stated bounds are wrong at boundary inputs that real call
sites do not reach. Both were verified against the shipping C.

\paragraph{OF-1.} The \code{montgomery\_reduce} comment states $-Q < r < Q$ (strict) over the inclusive
domain $-2^{31} Q \le a \le Q \cdot 2^{31}$. At $a = 2^{31} Q$ the function returns $r = Q$, violating the
strict upper bound. The guarantee that holds over the inclusive domain is $-Q \le r \le Q$; the strict
bound holds on the half-open domain, which is what we specify.

\paragraph{OF-2.} The \code{reduce32} comment states $-6283008 \le r \le 6283008$ for $a \le 2^{31} -
2^{22} - 1$. That precondition is one-sided, so $a = -2143289344$ is admissible and returns $r =
-6283009$, one below the stated bound. The stated bound holds only under the symmetric precondition
$|a| \le 2^{31} - 2^{22} - 1$.

Both comments are identical in \code{pq-crystals/dilithium} (the origin; PQClean re-namespaces it), so
the finding was reported there as a single issue. We did not file automatically; disclosure was
deliberate.

\section{Scope and limitations}\label{sec:scope}

The verified code is reference C and a Rust library, not the deployed AVX2/aarch64 or the
\code{mldsa-native} successor, where the genuinely error-prone optimized arithmetic lives. We prove
C/Rust $\equiv$ model and, for the C NTT, overflow-freedom; we do not prove NTT \emph{transform}
correctness against FIPS~204 Algorithms 41 and 42. The wide-domain Barrett reduction and the NTT-layer
transform proof are reported as a characterized boundary with escapes identified, not as completed
proofs; the C-side overflow-freedom statement is the composition argued in Section~\ref{sec:ntt}, not a
single mechanized theorem. The Rust verification rests on \code{mir-json} soundness and the three assumed
constant-time leaf specifications of Section~\ref{sec:rust}. We say nothing about constant-time behavior.

\section{Reproducibility}\label{sec:repro}

The development is public and all tool versions are pinned by setup scripts: SAW 1.5.1 with its bundled
Cryptol and \code{cryptol-to-isabelle}, Isabelle2025-2 with a pinned AFP snapshot, a fixed clang, and for
the Rust pipeline a pinned Rust nightly with \code{mir-json} at the schema commit SAW 1.5.1 bundles. The
target C is vendored with its commit hash; the Rust crate is vendored likewise. A single command,
\code{make verify}, runs the lift check, the SAW leg, and the Isabelle leg; the Rust proofs run from the
same toolchain. Continuous integration runs the fast SAW gates on every push and the full pipeline on
demand. The artifact is archived on Zenodo, DOI
\href{https://doi.org/10.5281/zenodo.20641411}{10.5281/zenodo.20641411}, and the source is at
\url{https://github.com/amarshat/pqc-assay}. The wide-domain Barrett timeout is reproduced by a single
self-contained script in the repository.

\section{Related work}\label{sec:related}

Apple's 2026 corecrypto verification~\cite{corecrypto} covers ML-KEM and ML-DSA in production code with
SAW, Cryptol, and Isabelle, and is the direct methodological precedent for the C pipeline; our work
applies the same public toolchain independently to reference C with a FIPS-derived specification, adds a
second toolchain over Rust MIR, and centers the tractability boundary rather than a single
implementation. Deployed implementations in OpenSSL, BoringSSL, and AWS-LC, and the \code{mldsa-native}
assembly, are verified by other toolchains such as CBMC and HOL-Light. A recent audit of reduction
placement in production ML-DSA~\cite{eprint1032} reports missing-reduction and overflow defects that
escaped test vectors, the defect class our overflow-freedom property targets. The difficulty of
discharging SMT obligations for nonlinear and wide modular arithmetic is well known in the
high-assurance-cryptography literature; our contribution is to characterize it concretely for ML-DSA's
reduction kernels, across two toolchains, with the structural cause and the escapes made explicit.

\section{Conclusion}\label{sec:conclusion}

Whether a machine-checked equivalence proof for ML-DSA's modular arithmetic discharges depends, sharply
and predictably, on the input domain. We verified the reduction layer of the C reference end to end, a
field, division, twiddle-table, and hint layer of the RustCrypto implementation, and an overflow-freedom
bound for the forward-NTT model; and we mapped the boundary where the transform layer ceases to be
tractable, identified its structural cause, and gave three proof restructurings that cross it. The same
boundary appears in both toolchains, which locates it in the goal rather than the tool. Mechanizing the
transform layer through the unbounded-integer and compositional escapes, and repointing the pipelines at
the deployed optimized arithmetic, are the natural next steps.

\paragraph{Acknowledgments.} We thank the Galois SAW and the Isabelle communities for guidance on the
toolchain.

\begin{thebibliography}{9}
\bibitem{fips204} National Institute of Standards and Technology.
  \emph{FIPS 204: Module-Lattice-Based Digital Signature Standard.} 2024.
\bibitem{corecrypto} Apple. \emph{Formal verification of corecrypto.}
  \url{https://github.com/apple/corecrypto/tree/2026-05}, 2026.
\bibitem{saw} Galois, Inc. \emph{SAW: The Software Analysis Workbench.}
  \url{https://github.com/GaloisInc/saw-script}.
\bibitem{mirjson} Galois, Inc. \emph{mir-json.} \url{https://github.com/GaloisInc/mir-json}.
\bibitem{pqclean} The PQClean project. \url{https://github.com/PQClean/PQClean}.
\bibitem{dilithium} CRYSTALS-Dilithium. \url{https://github.com/pq-crystals/dilithium}.
\bibitem{rustcrypto} The RustCrypto project. \emph{ml-dsa.}
  \url{https://github.com/RustCrypto/signatures}.
\bibitem{eprint1032} Sunwoo Lee, Hyuk Lim, and Seunghyun Yoon.
  \emph{When Removing Reductions Goes Wrong: Auditing Reduction Placement in Production ML-DSA
  Implementations.} IACR Cryptology ePrint Archive, Paper 2026/1032, 2026.
  \url{https://eprint.iacr.org/2026/1032}.
\end{thebibliography}

\end{document}
